% =============================================================================
% Manual P05 – Análisis de Imagen Digital y Fractografía SEM
% Técnicas de Medida – Ingeniería Aeronáutica
% Compilar con: latexmk -xelatex -interaction=nonstopmode manual_P05.tex
% =============================================================================
\documentclass[12pt,a4paper]{article}

% ---- Encoding & Language ----
\usepackage{fontspec}
\usepackage[spanish]{babel}
\addto\captionsspanish{\renewcommand{\tablename}{Tabla}}

% ---- Page Layout ----
\usepackage[top=2.5cm,bottom=2.5cm,left=2.5cm,right=2.5cm]{geometry}
\setlength{\headheight}{14.5pt}

% ---- Math & Physics ----
\usepackage{amsmath,amssymb,amsfonts}
\usepackage{mathtools}
\usepackage{siunitx}
\sisetup{output-decimal-marker={,}, group-separator={.}}

% ---- Graphics & Color ----
\usepackage{graphicx}
\usepackage[table]{xcolor}
\usepackage{float}
\usepackage{caption}
\usepackage{subcaption}

% ---- TikZ ----
\usepackage{tikz}
\usetikzlibrary{babel}
\usetikzlibrary{arrows.meta,positioning,calc,shapes.geometric,decorations.pathmorphing,
                decorations.markings,patterns,fit,backgrounds}

% ---- Tables ----
\usepackage{booktabs}
\usepackage{multirow}
\usepackage{array}
\usepackage{colortbl}
\usepackage{tabularx}
\usepackage{adjustbox}
\usepackage{longtable}

% ---- Code Listings ----
\usepackage{listings}
\lstset{
  basicstyle=\ttfamily\small,
  keywordstyle=\color{blue!70!black}\bfseries,
  commentstyle=\color{green!50!black},
  stringstyle=\color{red!60!black},
  backgroundcolor=\color{gray!10},
  frame=single,
  breaklines=true,
  numbers=left,
  numberstyle=\tiny\color{gray}
}

% ---- Hyperlinks ----
\usepackage[colorlinks=true,linkcolor=blue!70!black,citecolor=green!50!black,
            urlcolor=blue!50!black]{hyperref}

% ---- Bibliography ----
\usepackage{csquotes}
\usepackage[backend=biber,style=ieee,sorting=none]{biblatex}
\addbibresource{references_P05.bib}

% ---- Headers & Footers ----
\usepackage{fancyhdr}
\pagestyle{fancy}
\fancyhf{}
\fancyhead[L]{\small Técnicas de Medida -- Ingeniería Aeronáutica}
\fancyhead[R]{\small P05: Análisis de Imagen}
\fancyfoot[C]{\thepage}

% ---- Custom Colors ----
\definecolor{aeroblue}{RGB}{0,51,102}
\definecolor{aerored}{RGB}{153,0,0}
\definecolor{aerogreen}{RGB}{0,102,51}
\definecolor{aeroyellow}{RGB}{255,204,0}

% =============================================================================
\begin{document}

\begin{titlepage}
	\centering
	\vspace*{1cm}
	{\color{aeroblue}\rule{\linewidth}{2pt}}\\[0.4cm]
	{\Large\bfseries\color{aeroblue} TÉCNICAS DE MEDIDA}\\[0.3cm]
	{\color{aeroblue}\rule{\linewidth}{1pt}}\\[1.5cm]
	{\huge\bfseries P05 -- Análisis de Imagen Digital\\[0.3cm]
	y Fractografía SEM}\\[1cm]
	{\large Manual de Laboratorio}\\[0.5cm]
	{\large\itshape Ingeniería Aeronáutica}\\[2cm]
	% SEM + image processing schematic
	\begin{tikzpicture}[scale=0.85,>=Stealth]
		% SEM column
		\draw[fill=gray!20,draw=aeroblue,thick] (-0.6,3.0) rectangle (0.6,4.2);
		\node[aeroblue,font=\tiny\bfseries] at (0,3.6) {Cañón $e^-$};
		% Lens
		\draw[aeroblue,thick] (-1.0,2.4) arc (180:360:1.0 and 0.15);
		\draw[aeroblue,thick] (-1.0,2.4) arc (180:0:1.0 and 0.15);
		% Beam
		\draw[blue!50,line width=1.5pt] (0,3.0) -- (0,0.3);
		% Sample
		\draw[fill=orange!20,draw=black,thick] (-1.2,0) -- (1.2,0) -- (0.9,-0.4) -- (-0.9,-0.4) -- cycle;
		\node[font=\tiny] at (0,-0.7) {Espécimen de fractura};
		% Detector
		\draw[fill=aerogreen!20,draw=aerogreen,thick] (1.8,1.0) circle (0.35);
		\node[aerogreen,font=\tiny] at (1.8,0.5) {Detector};
		\draw[->,dashed,gray] (0.2,0.2) -- (1.5,0.8);
		% Arrow to processing
		\draw[aeroblue,ultra thick,->] (2.5,1.0) -- (4.0,1.0);
		\node[above,font=\tiny,aeroblue] at (3.25,1.1) {Imagen digital};
		% Processing pipeline boxes
		\draw[fill=aeroblue!10,draw=aeroblue,thick,rounded corners=3pt]
		      (4.0,0.2) rectangle (5.6,1.8);
		\node[aeroblue,font=\tiny\bfseries,align=center] at (4.8,1.0) {CLAHE\\Sobel\\Canny};
		\draw[aeroblue,thick,->] (5.6,1.0) -- (6.2,1.0);
		\draw[fill=aerored!10,draw=aerored,thick,rounded corners=3pt]
		      (6.2,0.2) rectangle (7.8,1.8);
		\node[aerored,font=\tiny\bfseries,align=center] at (7.0,1.0) {Otsu\\Morfología\\Etiquetado};
		% Output
		\draw[aerored,thick,->] (7.8,1.0) -- (8.5,1.0);
		\node[right,font=\tiny] at (8.5,1.0) {Regiones};
	\end{tikzpicture}\\[2cm]
	{\large Laboratorio de Técnicas de Medida}
	\vfill
	{\color{aeroblue}\rule{\linewidth}{2pt}}
\end{titlepage}

\tableofcontents
\newpage

% =============================================================================
\section{Introducción}
\label{sec:intro}

El análisis de superficies de fractura, o \textbf{fractografía}, mediante
Microscopía Electrónica de Barrido (SEM) es una herramienta indispensable
en la ingeniería aeronáutica para comprender los mecanismos de falla de los
materiales, evaluar la integridad estructural de componentes y guiar el
diseño de aleaciones más resistentes~\cite{asm_fractography,hull1999}.

Las micrografías SEM proporcionan imágenes de alta resolución y gran
profundidad de campo, revelando detalles topográficos que permiten
identificar el origen de la falla, el modo de propagación de la grieta y
las características microestructurales involucradas.  No obstante, la
información contenida en una micrografía es, en primera instancia,
\emph{cualitativa}.  Para extraer datos cuantitativos --- dimensiones de
hoyuelos, densidad de estrías, distribución de tamaño de grano --- es
necesario aplicar técnicas de \textbf{procesamiento digital de
imágenes}~\cite{gonzalez2018,russ2016}.

En esta práctica se combinan ambos enfoques: el estudiante procesará
micrografías SEM de superficies de fractura por fatiga utilizando un
\emph{pipeline} de análisis de imagen compuesto por ecualización adaptativa
(CLAHE), detección de bordes (Sobel, Canny), umbralización (Otsu),
operaciones morfológicas y etiquetado de regiones.  Adicionalmente, se
realizarán mediciones dimensionales interactivas utilizando las barras de
escala presentes en las imágenes.  El conjunto de técnicas permite simular
el proceso que un ingeniero de materiales realizaría para diagnosticar la
causa raíz de la falla de un componente aeronáutico.

% =============================================================================
\section{Objetivos de Aprendizaje}
\label{sec:objetivos}

Al finalizar esta práctica el estudiante será capaz de:

\begin{enumerate}
	\item Comprender la importancia de la microscopía electrónica de barrido
	      (SEM) para el análisis fractográfico en ingeniería aeronáutica.
	\item Interpretar metadatos de adquisición SEM: voltaje ($\mathrm{EHT}$),
	      distancia de trabajo ($\mathrm{WD}$), magnificación y detector.
	\item Aplicar la ecualización adaptativa de histograma (CLAHE) para
	      mejorar el contraste local en micrografías con iluminación desigual.
	\item Calcular y utilizar los operadores de detección de bordes de Sobel
	      y de Canny, comprendiendo sus fundamentos matemáticos.
	\item Aplicar la umbralización global de Otsu y justificar su validez
	      minimizando la varianza intra-clase.
	\item Realizar operaciones morfológicas (apertura, cierre, erosión,
	      dilatación) para refinar máscaras binarias.
	\item Etiquetar regiones conectadas y extraer descriptores de forma
	      (área, perímetro, excentricidad, eje mayor).
	\item Calibrar la escala espacial (píxeles/\si{\micro\m}) y medir
	      dimensiones de características microestructurales.
	\item Diagnosticar modos de fractura (dúctil, frágil, fatiga) basándose
	      en evidencia topográfica y datos cuantitativos.
	\item Evaluar los riesgos asociados a la operación del SEM y al manejo
	      de muestras.
\end{enumerate}

% =============================================================================
\section{Fundamentos Teóricos}
\label{sec:teoria}

\subsection{Microscopía Electrónica de Barrido (SEM)}

A diferencia de la microscopía óptica, el SEM utiliza un haz de electrones
focalizado para barrer la superficie de la muestra.  La interacción del haz
con los átomos de la superficie genera señales --- electrones secundarios
(SE), electrones retrodispersados (BSE), rayos~X característicos --- que
son detectadas y transformadas en una imagen
digital~\cite{goldstein2018}.

\begin{figure}[H]
	\centering
	\begin{tikzpicture}[scale=0.80,>=Stealth]
		% Column
		\draw[fill=gray!15,draw=black,thick] (-0.8,4.5) rectangle (0.8,5.5);
		\node[font=\small] at (0,5.0) {Cañón $e^-$};
		% Condenser lens
		\draw[thick,aeroblue] (-1.2,3.8) arc (180:360:1.2 and 0.18);
		\draw[thick,aeroblue] (-1.2,3.8) arc (180:0:1.2 and 0.18);
		\node[right,font=\tiny,aeroblue] at (1.3,3.8) {Lente condensadora};
		% Objective lens
		\draw[thick,aeroblue] (-1.0,2.2) arc (180:360:1.0 and 0.15);
		\draw[thick,aeroblue] (-1.0,2.2) arc (180:0:1.0 and 0.15);
		\node[right,font=\tiny,aeroblue] at (1.1,2.2) {Lente objetivo};
		% Scan coils
		\draw[fill=aeroyellow!30,draw=black] (-0.5,3.0) rectangle (-0.2,3.4);
		\draw[fill=aeroyellow!30,draw=black] ( 0.2,3.0) rectangle ( 0.5,3.4);
		\node[font=\tiny] at (0,2.7) {Bobinas de barrido};
		% Beam
		\draw[blue!60,line width=1.5pt] (0,4.5) -- (0,0.3);
		% Sample
		\draw[fill=orange!15,draw=black,thick]
			(-1.5,0) -- (1.5,0) -- (1.2,-0.5) -- (-1.2,-0.5) -- cycle;
		\node[font=\small] at (0,-0.8) {Muestra};
		% SE detector
		\draw[fill=aerogreen!20,draw=aerogreen,thick] (2.5,1.2) ellipse (0.4 and 0.3);
		\node[aerogreen,font=\tiny,below] at (2.5,0.85) {Detector SE};
		\draw[->,dashed,aerogreen] (0.3,0.2) -- (2.1,1.0);
		% BSE arrow
		\draw[->,dashed,aerored] (0.1,0.2) -- (-2.0,1.5);
		\node[aerored,font=\tiny] at (-2.5,1.5) {BSE};
		% X-ray arrow
		\draw[->,dashed,blue!50] (-0.1,0.2) -- (-2.2,0.5);
		\node[blue!50,font=\tiny] at (-2.6,0.5) {Rayos~X};
	\end{tikzpicture}
	\caption{Esquema simplificado de un microscopio electrónico de barrido.
		El haz de electrones barre la muestra; las señales SE, BSE y rayos~X
		se detectan para formar la imagen.}
	\label{fig:sem_esquema}
\end{figure}

Las ventajas del SEM para fractografía incluyen~\cite{goldstein2018}:

\begin{itemize}
	\item \textbf{Alta resolución:} hasta $\sim\!\SI{1}{\nano\m}$ con
	      detectores de campo cercano (FE-SEM).
	\item \textbf{Gran profundidad de campo:} produce imágenes con apariencia
	      tridimensional, ideal para superficies de fractura irregulares.
	\item \textbf{Análisis composicional:} combinado con EDS/EDX permite
	      identificar elementos químicos en inclusiones o precipitados.
	\item \textbf{Amplio rango de magnificación:} desde $\sim\!10\times$
	      hasta $>300\,000\times$.
\end{itemize}

\subsubsection{Parámetros de adquisición}

Cada micrografía SEM incluye metadatos de adquisición relevantes para la
interpretación:

\begin{table}[H]
	\centering
	\caption{Parámetros de adquisición SEM y su influencia en la imagen.}
	\label{tab:parametros_sem}
	\begin{tabularx}{\textwidth}{l l X}
		\toprule
		\textbf{Parámetro} & \textbf{Unidad}           & \textbf{Efecto}              \\
		\midrule
		EHT (voltaje)      & \si{\kilo\volt}           & Mayor EHT $\to$ mayor
		penetración, mejor resolución, mayor daño por radiación.                      \\
		WD (dist.\ trabajo) & \si{\mm}                  & Menor WD $\to$ mejor
		resolución; mayor WD $\to$ mayor profundidad de campo.                        \\
		Mag (magnificación)  & $\times$                  & Define el campo de visión
		y la resolución espacial efectiva.                                             \\
		Signal A (detector)  & SE / BSE                  & SE: topografía; BSE:
		contraste composicional (número atómico).                                     \\
		Escala (scale bar)   & \si{\micro\m}             & Relación entre $N$ píxeles
		y la longitud real, fundamental para metrología.                               \\
		\bottomrule
	\end{tabularx}
\end{table}

\subsection{Mecanismos de Fractura en Materiales Aeroespaciales}

La fractografía permite clasificar el modo de falla de un componente
estructural.  Los tres mecanismos principales son~\cite{asm_fractography,
anderson2017,dowling2013}:

\subsubsection{Fractura dúctil}

Se produce con deformación plástica significativa.  A nivel microscópico se
manifiesta por \textbf{hoyuelos} (\emph{dimples}): cavidades formadas por
la nucleación, crecimiento y coalescencia de micro-vacíos en torno a
inclusiones o precipitados.

\begin{equation}
	\text{Energía absorbida:}\quad G_{\mathrm{dúctil}} \gg G_{\mathrm{frágil}}
	\label{eq:energia_fractura}
\end{equation}

La forma de los hoyuelos contiene información sobre el estado de
tensiones~\cite{hull1999}:

\begin{itemize}
	\item \textbf{Equiaxiales}: tracción uniaxial pura.
	\item \textbf{Alargados/parabólicos}: estado de cortante o tracción
	      oblicua.
	\item \textbf{Pequeños y densos}: alta ductilidad del material base.
\end{itemize}

\subsubsection{Fractura frágil}

Ocurre con poca o ninguna deformación plástica.  Las superficies suelen
ser relativamente planas~\cite{anderson2017}:

\begin{itemize}
	\item \textbf{Clivaje transgranular}: la grieta sigue planos
	      cristalográficos de baja energía superficial, generando \emph{facetas
	      planas} con \emph{marcas de río} que convergen hacia el origen.
	\item \textbf{Fractura intergranular}: la grieta se propaga a lo largo de
	      los límites de grano.  Indica fragilización por segregación de
	      impurezas, hidrógeno, o tratamiento térmico inadecuado.
\end{itemize}

\subsubsection{Fractura por fatiga}

Es la falla progresiva bajo cargas cíclicas, responsable de la
mayoría de las fallas en servicio de componentes
aeronáuticos~\cite{suresh1998,dowling2013}:

\begin{equation}
	\frac{da}{dN} = C\,(\Delta K)^m \qquad\text{(Ley de Paris)}
	\label{eq:paris}
\end{equation}

donde $a$ es la longitud de grieta, $N$ el número de ciclos, $\Delta K$ el
rango del factor de intensidad de tensiones, y $C$, $m$ constantes del
material.  Las superficies de fractura por fatiga muestran:

\begin{itemize}
	\item \textbf{Estrías de fatiga}: marcas finas y paralelas; cada estría
	      corresponde a un ciclo de carga.  El espaciado entre estrías
	      ($\approx da/dN$) aumenta con $\Delta K$.
	\item \textbf{Marcas de playa} (\emph{beach marks}): bandas
	      macroscópicas concéntricas que registran cambios en la amplitud de
	      carga o interrupciones del servicio.
	\item \textbf{Origen de la grieta}: zona frecuentemente asociada a
	      defectos superficiales, concentradores de tensión o inclusiones.
\end{itemize}

\begin{figure}[H]
	\centering
	\begin{tikzpicture}[scale=0.85,>=Stealth]
		% Ductile fracture
		\begin{scope}[shift={(0,0)}]
			\draw[thick] (0,0) rectangle (3.2,2.5);
			\node[above,font=\small\bfseries] at (1.6,2.5) {Fractura Dúctil};
			% Dimples
			\foreach \x/\y/\r in {0.5/1.2/0.25, 1.0/0.8/0.3, 1.5/1.5/0.2,
				2.0/1.0/0.35, 2.5/1.8/0.2, 1.2/1.9/0.15, 2.2/0.5/0.25,
				0.7/0.4/0.2, 1.8/1.3/0.18, 2.7/1.2/0.22} {
				\draw[gray,fill=gray!20] (\x,\y) ellipse ({\r} and {\r*0.7});
			}
			\node[font=\tiny,aeroblue] at (1.6,0.1) {Hoyuelos (dimples)};
		\end{scope}
		% Brittle fracture
		\begin{scope}[shift={(4.2,0)}]
			\draw[thick] (0,0) rectangle (3.2,2.5);
			\node[above,font=\small\bfseries] at (1.6,2.5) {Fractura Frágil};
			% Cleavage facets
			\draw[gray] (0.3,0.2) -- (1.0,1.8) -- (1.5,0.5) -- cycle;
			\draw[gray] (1.0,1.8) -- (2.2,2.0) -- (1.5,0.5) -- cycle;
			\draw[gray] (1.5,0.5) -- (2.2,2.0) -- (2.8,0.8) -- cycle;
			\draw[gray] (2.2,2.0) -- (2.9,2.3) -- (2.8,0.8) -- cycle;
			\draw[gray] (0.3,0.2) -- (1.5,0.5) -- (0.5,2.2) -- cycle;
			% River marks
			\draw[aerored,thin,->] (0.8,1.0) -- (1.2,1.2);
			\draw[aerored,thin,->] (1.3,0.8) -- (1.7,1.0);
			\draw[aerored,thin,->] (2.0,1.2) -- (2.4,1.5);
			\node[font=\tiny,aerored] at (1.6,0.1) {Facetas + marcas de río};
		\end{scope}
		% Fatigue
		\begin{scope}[shift={(8.4,0)}]
			\draw[thick] (0,0) rectangle (3.2,2.5);
			\node[above,font=\small\bfseries] at (1.6,2.5) {Fatiga};
			% Striations
			\foreach \y in {0.3,0.55,0.8,1.05,1.3,1.55,1.8,2.05,2.3} {
				\draw[gray,thin] (0.2,\y) .. controls (1.0,{\y+0.15}) and
					(2.2,{\y-0.1}) .. (3.0,\y);
			}
			% Origin marker
			\fill[aerored] (0.2,1.3) circle (0.08);
			\draw[aerored,->] (0.2,1.3) -- (0.8,1.3);
			\node[font=\tiny,aerogreen] at (1.6,0.1) {Estrías de fatiga};
		\end{scope}
	\end{tikzpicture}
	\caption{Representación esquemática de los tres mecanismos de fractura
		principales observados en micrografías SEM.}
	\label{fig:modos_fractura}
\end{figure}

\subsection{Conceptos de Imagen Digital}

Una imagen digital en escala de grises es una función
$I : \{0,\ldots,M-1\} \times \{0,\ldots,N-1\} \to [0,L-1]$, donde
$M \times N$ es la resolución en píxeles y $L$ el número de niveles de
gris (típicamente $L = 256$ para 8~bits)~\cite{gonzalez2018}.

\subsubsection{Histograma}

El histograma $h(k)$ cuenta el número de píxeles con intensidad $k$:

\begin{equation}
	h(k) = \#\{(x,y) : I(x,y) = k\}, \quad k = 0, 1, \ldots, L-1
	\label{eq:histograma}
\end{equation}

Su forma revela el contraste de la imagen: un histograma estrecho indica
bajo contraste; uno bimodal sugiere dos poblaciones (fondo y objeto)
separables por umbralización.

\subsection{Pipeline de Procesamiento de Imagen}

El análisis automatizado implementado en el script sigue un \emph{pipeline}
secuencial de seis etapas (Figura~\ref{fig:pipeline}).  Cada etapa
transforma la imagen para extraer progresivamente la información
morfológica relevante.

\begin{figure}[H]
	\centering
	\begin{tikzpicture}[node distance=1.3cm,>=Stealth,
		block/.style={rectangle,draw=aeroblue,thick,fill=aeroblue!8,
			text width=4.5cm,align=center,minimum height=0.8cm,
			font=\small,rounded corners=3pt},
		result/.style={rectangle,draw=aerored,thick,fill=aerored!8,
			text width=4.5cm,align=center,minimum height=0.8cm,
			font=\small,rounded corners=3pt}]
		\node[block] (load)   {1. Carga en gris + normalización};
		\node[block,below=of load]  (clahe) {2. CLAHE (ecualización adaptativa)};
		\node[block,below=of clahe] (edges)
			{3. Detección de bordes (Sobel / Canny)};
		\node[block,below=of edges] (otsu)  {4. Umbralización de Otsu};
		\node[block,below=of otsu]  (morph) {5. Morfología (Opening, disk=2)};
		\node[result,below=of morph] (label)
			{6. Etiquetado de regiones conectadas};
		% Arrows
		\foreach \a/\b in {load/clahe,clahe/edges,edges/otsu,otsu/morph,morph/label} {
			\draw[->,aeroblue,thick] (\a) -- (\b);
		}
		% Side annotations
		\node[right=1.5cm of load,font=\tiny,text width=3cm]
			{$I(x,y)\in[0,1]$ float};
		\node[right=1.5cm of clahe,font=\tiny,text width=3cm]
			{Contraste local mejorado};
		\node[right=1.5cm of edges,font=\tiny,text width=3cm]
			{Mapa de gradientes / bordes binarios};
		\node[right=1.5cm of otsu,font=\tiny,text width=3cm]
			{Máscara binaria (objeto/fondo)};
		\node[right=1.5cm of morph,font=\tiny,text width=3cm]
			{Ruido eliminado, regiones limpias};
		\node[right=1.5cm of label,font=\tiny,text width=3cm,aerored]
			{$N$ regiones con descriptores de forma};
	\end{tikzpicture}
	\caption{Pipeline de procesamiento de imagen implementado en el script
		\texttt{P05\_Analisis\_Imagen.py}.}
	\label{fig:pipeline}
\end{figure}

\subsection{Ecualización Adaptativa: CLAHE}
\label{subsec:clahe}

La ecualización clásica de histograma redistribuye los niveles de gris
para maximizar el contraste global.  Sin embargo, en micrografías SEM con
iluminación desigual (e.g., sombras por topografía), puede sobre-amplificar
regiones ya brillantes.

La \textbf{CLAHE} (\emph{Contrast Limited Adaptive Histogram
Equalization})~\cite{zuiderveld1994} resuelve este problema dividiendo la
imagen en regiones locales (\emph{tiles}) y ecualizando cada una de forma
independiente, con un límite de contraste (\texttt{clip\_limit}) que evita
la amplificación excesiva del ruido:

\begin{equation}
	\boxed{I_{\mathrm{CLAHE}}(x,y)
		= T_{\mathrm{tile}(x,y)}\!\bigl[I(x,y)\bigr]}
	\label{eq:clahe}
\end{equation}

donde $T_{\mathrm{tile}}$ es la función de transferencia de ecualización
calculada para el \emph{tile} local, con redistribución del exceso de
contraste por encima de \texttt{clip\_limit}.  Transiciones entre
\emph{tiles} se suavizan mediante interpolación bilineal.

En esta práctica se utiliza \texttt{clip\_limit}~$= 0{,}03$ (normalizado):

\begin{lstlisting}[language=Python,caption={Aplicación de CLAHE con
  scikit-image.}]
from skimage import exposure
img_clahe = exposure.equalize_adapthist(
    img_gray, clip_limit=0.03)
\end{lstlisting}

\subsection{Operador de Sobel}
\label{subsec:sobel}

El filtro de Sobel~\cite{sobel1968} estima el gradiente de la imagen
mediante la convolución con dos kernels $3\times3$ que calculan las
derivadas parciales en las direcciones $x$ e~$y$:

\begin{equation}
	G_x = \begin{bmatrix} -1 & 0 & +1 \\ -2 & 0 & +2 \\ -1 & 0 & +1 \end{bmatrix}
	\ast I, \qquad
	G_y = \begin{bmatrix} +1 & +2 & +1 \\ 0 & 0 & 0 \\ -1 & -2 & -1 \end{bmatrix}
	\ast I
	\label{eq:sobel_kernels}
\end{equation}

La magnitud del gradiente combina ambas componentes:

\begin{equation}
	\boxed{|\nabla I| = \sqrt{G_x^2 + G_y^2}}
	\label{eq:sobel_mag}
\end{equation}

Valores altos de $|\nabla I|$ corresponden a transiciones abruptas de
intensidad (bordes).  Sobel es sensible al ruido; por ello se aplica
sobre la imagen CLAHE (ya con contraste mejorado).

\subsection{Detector de Bordes de Canny}
\label{subsec:canny}

El algoritmo de Canny~\cite{canny1986} es un detector de bordes óptimo
que busca maximizar la detección, minimizar las falsas respuestas y
producir bordes de un solo píxel de ancho.  Consta de cuatro pasos:

\begin{enumerate}
	\item \textbf{Suavizado Gaussiano:} eliminación de ruido con un kernel
	      $G_\sigma$ de desviación estándar $\sigma$.
	      \begin{equation}
	      	G_\sigma(x,y) = \frac{1}{2\pi\sigma^2}
	      	\exp\!\left(-\frac{x^2 + y^2}{2\sigma^2}\right)
	      	\label{eq:gauss_kernel}
	      \end{equation}
	\item \textbf{Gradiente:} cálculo de $|\nabla I|$ y orientación
	      $\theta = \arctan(G_y/G_x)$.
	\item \textbf{Supresión no máxima:} se conservan solo los píxeles cuyo
	      gradiente es máximo local en la dirección de $\theta$.
	\item \textbf{Histéresis con doble umbral:} un umbral alto
	      $T_{\mathrm{high}}$ confirma bordes fuertes; un umbral bajo
	      $T_{\mathrm{low}}$ extiende los bordes conectados.
\end{enumerate}

En esta práctica se emplea $\sigma = 2{,}0$, que suaviza suficiente para
micrografías SEM con textura granular fina.

\subsection{Umbralización de Otsu}
\label{subsec:otsu}

El método de Otsu~\cite{otsu1979} calcula automáticamente el umbral
$t^*$ que separa la imagen en dos clases (fondo $C_0$ y objeto $C_1$)
\textbf{minimizando la varianza intra-clase} (o, equivalentemente,
\textbf{maximizando la varianza inter-clase}).

\paragraph{Formulación matemática.}  Sea $p_i = h(i)/N_\mathrm{total}$
la probabilidad del nivel de gris~$i$.  Para un umbral~$t$:

\begin{align}
	\omega_0(t) & = \sum_{i=0}^{t} p_i, &
	\omega_1(t) & = \sum_{i=t+1}^{L-1} p_i
	\label{eq:otsu_omega}
\end{align}

\begin{align}
	\mu_0(t) & = \frac{1}{\omega_0}\sum_{i=0}^{t} i\,p_i, &
	\mu_1(t) & = \frac{1}{\omega_1}\sum_{i=t+1}^{L-1} i\,p_i
	\label{eq:otsu_mu}
\end{align}

La varianza inter-clase es:

\begin{equation}
	\sigma_B^2(t) = \omega_0(t)\,\omega_1(t)\,
	\bigl[\mu_0(t) - \mu_1(t)\bigr]^2
	\label{eq:otsu_sigma}
\end{equation}

El umbral óptimo es:

\begin{equation}
	\boxed{t^* = \arg\max_{0 \le t < L}\;\sigma_B^2(t)}
	\label{eq:otsu_opt}
\end{equation}

La imagen binarizada se obtiene como
$B(x,y) = \mathbf{1}[I(x,y) > t^*]$.

\subsection{Operaciones Morfológicas}
\label{subsec:morfologia}

Las operaciones morfológicas procesan imágenes binarias utilizando un
\textbf{elemento estructurante} (\emph{structuring element}, SE); en
esta práctica, un disco de radio~2 (\texttt{disk(2)}).  Las operaciones
básicas son~\cite{soille2004}:

\paragraph{Erosión:}
\begin{equation}
	(I \ominus SE)(x,y)
	= \min_{(s,t)\in SE}\; I(x+s,\,y+t)
	\label{eq:erosion}
\end{equation}

\paragraph{Dilatación:}
\begin{equation}
	(I \oplus SE)(x,y)
	= \max_{(s,t)\in SE}\; I(x+s,\,y+t)
	\label{eq:dilatacion}
\end{equation}

\paragraph{Apertura (Opening):} erosión seguida de dilatación.  Elimina
pequeños objetos (ruido) preservando la forma de los objetos grandes:

\begin{equation}
	\boxed{I \circ SE = (I \ominus SE) \oplus SE}
	\label{eq:opening}
\end{equation}

\paragraph{Cierre (Closing):} dilatación seguida de erosión.  Rellena
pequeños agujeros en regiones:

\begin{equation}
	I \bullet SE = (I \oplus SE) \ominus SE
	\label{eq:closing}
\end{equation}

\subsection{Etiquetado de Regiones Conectadas}
\label{subsec:etiquetado}

Sobre la imagen binarizada y limpiada morfológicamente, se aplica el
\textbf{etiquetado de componentes conexos}: cada grupo de píxeles
``1'' que comparte vecindad (4- u 8-conectividad) se asigna a una
etiqueta única $\ell \in \{1, 2, \ldots, N_\mathrm{reg}\}$.

Para cada región etiquetada se calculan propiedades geométricas
(descriptores de forma)~\cite{russ2016}:

\begin{table}[H]
	\centering
	\caption{Descriptores de forma calculados por \texttt{regionprops}.}
	\label{tab:descriptores}
	\begin{tabularx}{\textwidth}{l l X}
		\toprule
		\textbf{Descriptor}          & \textbf{Unidades}     & \textbf{Interpretación}                \\
		\midrule
		Área ($A_r$)                 & px$^2$                & Tamaño de la región; proporcional
		al área real mediante la escala.                                                              \\
		Perímetro ($P_r$)            & px                    & Longitud del borde de la región.       \\
		Excentricidad ($e$)          & $[0,1)$               & $0$: circular; cercano a $1$:
		muy alargada.                                                                                 \\
		Eje mayor ($d_\mathrm{max}$) & px                    & Longitud del eje mayor de la elipse
		equivalente; indicador del tamaño máximo.                                                     \\
		\bottomrule
	\end{tabularx}
\end{table}

\subsection{Calibración de Escala}
\label{subsec:calibracion}

Las barras de escala impresas en las micrografías SEM establecen la
relación entre píxeles y longitud real.  El procedimiento de calibración
es:

\begin{enumerate}
	\item Hacer clic en los dos extremos de la barra de escala (puntos
	      $P_1$ y $P_2$).
	\item Calcular la distancia en píxeles:
	      \begin{equation}
	      	d_\mathrm{px}
	      	= \sqrt{(x_2 - x_1)^2 + (y_2 - y_1)^2}
	      	\label{eq:dist_px}
	      \end{equation}
	\item Ingresar la longitud real $d_\mathrm{real}$ (en \si{\micro\m})
	      indicada por la barra.
	\item Calcular la relación de escala:
	      \begin{equation}
	      	\boxed{s = \frac{d_\mathrm{px}}{d_\mathrm{real}}
	      	\quad\left[\frac{\mathrm{px}}{\si{\micro\m}}\right]}
	      	\label{eq:escala}
	      \end{equation}
\end{enumerate}

Toda medición posterior se convierte a unidades reales dividiendo la
distancia en píxeles por~$s$.  Un error de $\pm n$ píxeles en la
selección de la barra produce un error relativo:

\begin{equation}
	\varepsilon_s = \frac{2n}{d_\mathrm{px}} \times 100\%
	\label{eq:error_escala}
\end{equation}

% =============================================================================
\section{Escenarios y Datos de la Práctica}
\label{sec:escenarios}

\subsection{Conjunto de datos}

Se dispone de cinco micrografías SEM de superficies de fractura por fatiga
almacenadas en \texttt{files/p5/}:

\begin{table}[H]
	\centering
	\caption{Micrografías SEM disponibles para el análisis.}
	\label{tab:imagenes}
	\begin{tabular}{l r l l}
		\toprule
		\textbf{Archivo}    & \textbf{Tamaño} & \textbf{Resolución}
		                    & \textbf{Observaciones}                              \\
		\midrule
		\texttt{fatigue3.png} & \SI{195}{\kilo\byte}
		                    & Moderada                                            & Estrías bien definidas \\
		\texttt{fatigue4.png} & \SI{68}{\kilo\byte}
		                    & Baja                                                & Buen contraste         \\
		\texttt{fatigue5.png} & \SI{2.3}{\mega\byte}
		                    & Alta                                                & Imagen de gran tamaño  \\
		\texttt{fatigue6.png} & \SI{1.9}{\mega\byte}
		                    & Alta                                                & Superficie compleja    \\
		\texttt{fatigue7.png} & \SI{964}{\kilo\byte}
		                    & Media-alta                                          & Textura fina           \\
		\bottomrule
	\end{tabular}
\end{table}

Cada micrografía muestra diferentes zonas o condiciones de una superficie
de fractura por fatiga.  Las diferencias de tamaño reflejan distintas
magnificaciones y condiciones de adquisición.

\subsection{Salidas del análisis}

El script genera, para cada imagen, 6 archivos en
\texttt{Practicas/P05\_Analisis\_Imagen/data/images/}:

\begin{table}[H]
	\centering
	\caption{Ficheros de salida generados por imagen procesada.}
	\label{tab:salidas}
	\small
	\begin{tabular}{l l}
		\toprule
		\textbf{Patrón de nombre}                     & \textbf{Contenido}               \\
		\midrule
		\texttt{\{nombre\}\_00\_resumen.png}          & Figura compuesta 3$\times$2
		(histograma, CLAHE, Sobel, Canny, Otsu, Opening)                                 \\
		\texttt{\{nombre\}\_01\_clahe.png}             & Imagen tras CLAHE                \\
		\texttt{\{nombre\}\_02\_sobel.png}             & Mapa del filtro Sobel            \\
		\texttt{\{nombre\}\_03\_canny.png}             & Bordes detectados por Canny      \\
		\texttt{\{nombre\}\_04\_otsu.png}              & Binarización de Otsu             \\
		\texttt{\{nombre\}\_05\_opening.png}           & Máscara morfológica (apertura)   \\
		\bottomrule
	\end{tabular}
\end{table}

En total, el procesamiento de las 5 imágenes genera \textbf{30 ficheros}
de resultados.

% =============================================================================
\section{Metodología de Cálculo}
\label{sec:metodologia}

\subsection{Paso 1: Carga y normalización}

La imagen se carga en escala de grises y se normaliza al rango $[0,1]$
(tipo \texttt{float64}):

\begin{equation}
	I_\mathrm{norm}(x,y) = \frac{I_\mathrm{raw}(x,y)}{L - 1}
	\label{eq:normalizacion}
\end{equation}

\subsection{Paso 2: Mejora de contraste (CLAHE)}

Se aplica la ecualización adaptativa (§\ref{subsec:clahe}) con
\texttt{clip\_limit}~$= 0{,}03$.  La imagen resultante se convierte a
\texttt{uint8} [$0,255$] para los pasos de detección de bordes, que
esperan niveles de gris enteros:

\begin{equation}
	I_\mathrm{u8}(x,y) = \mathrm{round}\bigl(255 \cdot I_\mathrm{CLAHE}(x,y)\bigr)
	\label{eq:uint8}
\end{equation}

\subsection{Paso 3: Detección de bordes}

Se aplican dos detectores complementarios sobre $I_\mathrm{u8}$:

\begin{itemize}
	\item \textbf{Sobel} (ec.~\ref{eq:sobel_mag}): produce un mapa
	      continuo de intensidad de bordes.  Útil para visualización
	      cualitativa.
	\item \textbf{Canny} (§\ref{subsec:canny}, $\sigma = 2{,}0$): produce
	      una imagen binaria con bordes de un píxel de ancho.  Adecuado
	      para la detección precisa de contornos de estrías.
\end{itemize}

\subsection{Paso 4: Segmentación (Otsu)}

La umbralización de Otsu (ec.~\ref{eq:otsu_opt}) segmenta la imagen
CLAHE en \emph{objeto} y \emph{fondo}, generando una máscara binaria
$B(x,y) \in \{0, 1\}$.

\subsection{Paso 5: Limpieza morfológica (Opening)}

La apertura con disco de radio~2 (ec.~\ref{eq:opening}) elimina regiones
espurias de menos de $\sim\!\pi(2)^2 \approx 12{,}6$~px$^2$, manteniendo
las regiones significativas correspondientes a características reales de
la superficie de fractura.

\subsection{Paso 6: Etiquetado y cuantificación}

El algoritmo de etiquetado de componentes conexos asigna una etiqueta
única a cada región.  Los descriptores de forma (Tabla~\ref{tab:descriptores})
se calculan automáticamente mediante \texttt{regionprops\_table} de
scikit-image~\cite{scikit_image}.

\subsection{Diagrama de flujo del procesamiento}

\begin{figure}[H]
	\centering
	\begin{tikzpicture}[node distance=1.2cm,>=Stealth,
		block/.style={rectangle,draw=aeroblue,thick,fill=aeroblue!8,
			text width=5.8cm,align=center,minimum height=0.85cm,
			font=\small,rounded corners=3pt}]
		\node[block] (glob) {Buscar imágenes: \texttt{glob(*.png)}};
		\node[block,below=of glob] (load) {Cargar en gris + normalizar $[0,1]$};
		\node[block,below=of load] (clahe) {CLAHE (clip\_limit\,=\,0,03)};
		\node[block,below=of clahe] (u8) {Convertir a uint8 $[0,255]$};
		\node[block,below=of u8] (sobel) {Sobel $\to$ guardar \_02};
		\node[block,below=of sobel] (canny) {Canny ($\sigma=2$) $\to$ guardar \_03};
		\node[block,below=of canny] (otsu) {Otsu $\to$ guardar \_04};
		\node[block,below=of otsu] (morph) {Opening (disk=2) $\to$ guardar \_05};
		\node[block,below=of morph] (label) {Etiquetado $\to$ $N$ regiones};
		\node[block,below=of label] (comp) {Construir figura compuesta \_00};
		% Arrows
		\foreach \a/\b in {glob/load,load/clahe,clahe/u8,u8/sobel,
			sobel/canny,canny/otsu,otsu/morph,morph/label,label/comp} {
			\draw[->,aeroblue,thick] (\a) -- (\b);
		}
		% Loop
		\draw[aerored,thick,->] (comp.east) -- ++(1.5,0) |- (load.east)
			node[pos=0.25,right,font=\tiny,aerored] {siguiente imagen};
	\end{tikzpicture}
	\caption{Diagrama de flujo detallado del script de análisis de imagen.}
	\label{fig:diagrama_detallado}
\end{figure}

% =============================================================================
\section{Ejemplo Resuelto: Análisis de \texttt{fatigue3.png}}
\label{sec:ejemplo}

\subsection{Ejecución del pipeline}

Se ejecuta el script desde la raíz del repositorio:

\begin{lstlisting}[language=bash]
python Practicas/P05_Analisis_Imagen/src/P05_Analisis_Imagen.py
\end{lstlisting}

Para \texttt{fatigue3.png}, la salida del etiquetado reporta:

\begin{verbatim}
--- Análisis de Regiones (sobre Otsu+Opening) ---
Número de objetos encontrados: 418
\end{verbatim}

\subsection{Interpretación de las etapas}

\begin{enumerate}
	\item \textbf{Histograma original}: la distribución de intensidades se
	      concentra en el rango medio-bajo, indicando una imagen relativamente
	      oscura con contraste limitado.
	\item \textbf{CLAHE}: mejora significativamente la visibilidad de
	      detalles finos (estrías, bordes de hoyuelos) sin sobre-saturar las
	      zonas brillantes.
	\item \textbf{Sobel}: revela la textura de la superficie; las zonas con
	      mayor gradiente corresponden a transiciones entre facetas y bordes
	      de estrías.
	\item \textbf{Canny ($\sigma = 2{,}0$)}: produce contornos nítidos de
	      un solo píxel, delineando los bordes de las características
	      fractográficas principales.
	\item \textbf{Otsu}: segmenta la imagen en regiones claras (elevaciones)
	      y oscuras (surcos/hoyuelos).  El umbral óptimo se calcula
	      automáticamente.
	\item \textbf{Opening (disk=2)}: elimina el ruido de píxeles aislados,
	      dejando \textbf{418 regiones} conectadas que corresponden a
	      características morfológicas reales.
\end{enumerate}

\subsection{Resultados comparativos}

\begin{table}[H]
	\centering
	\caption{Número de regiones detectadas para cada micrografía.}
	\label{tab:regiones}
	\begin{tabular}{l S[table-format=4.0] l}
		\toprule
		\textbf{Imagen}     & {$N_\mathrm{regiones}$} & \textbf{Observación}          \\
		\midrule
		fatigue3.png        & 418                     & Moderada densidad de rasgos   \\
		fatigue4.png        & 153                     & Baja densidad; buen contraste \\
		fatigue5.png        & 2066                    & Alta resolución, muy detallada \\
		fatigue6.png        & 2897                    & Superficie extremadamente compleja \\
		fatigue7.png        & 549                     & Textura fina, densidad media     \\
		\bottomrule
	\end{tabular}
\end{table}

\textbf{Análisis:} Las imágenes de alta resolución (\texttt{fatigue5},
\texttt{fatigue6}) producen un número significativamente mayor de regiones
debido a que el algoritmo resuelve más detalles topográficos.  Esto
ilustra la dependencia de las métricas cuantitativas con la magnificación
y la resolución de adquisición --- un punto fundamental en metrología de
imágenes~\cite{russ2016}.

\subsection{Evaluación de Riesgo de los Escenarios Base}
\label{sec:riesgo_base}

El entorno de trabajo con SEM y muestras metalográficas presenta riesgos
específicos que deben evaluarse.

\begin{table}[H]
	\centering
	\caption{Evaluación de riesgo para los escenarios de la práctica.}
	\label{tab:riesgo_base}
	\small
	\begin{tabularx}{\textwidth}{l c c c c X}
		\toprule
		\textbf{Escenario}     & \textbf{P} & \textbf{S} & \textbf{$R$}
		                       & \textbf{Nivel}
		                       & \textbf{Controles}                                        \\
		\midrule
		\rowcolor{aeroyellow!15}
		Operación del SEM (alto voltaje)
		                       & 2          & 4          & 8
		                       & \textcolor{aeroyellow!80!black}{\textbf{Medio}}
		                       & Formación, interlock de cámara, puesta a tierra            \\
		\rowcolor{aerogreen!10}
		Preparación de muestras (pulidora)
		                       & 2          & 3          & 6
		                       & \textcolor{aeroyellow!80!black}{\textbf{Medio}}
		                       & Gafas de seguridad, guantes, SOPs                         \\
		\rowcolor{aerogreen!10}
		Manipulación de reactivos (ataque metalográfico)
		                       & 2          & 3          & 6
		                       & \textcolor{aeroyellow!80!black}{\textbf{Medio}}
		                       & Campana extractora, guantes, MSDS                         \\
		\rowcolor{aerogreen!10}
		Procesamiento digital (ergonomía)
		                       & 1          & 1          & 1
		                       & \textcolor{aerogreen}{\textbf{Bajo}}
		                       & Pausas, ergonomía de estación                              \\
		\bottomrule
	\end{tabularx}
\end{table}

% =============================================================================
\section{Casos de Estudio por Equipos}
\label{sec:equipos}

Cada equipo de trabajo desarrollará uno de los siguientes escenarios
utilizando el script y/o el notebook de análisis.  Deben:

\begin{enumerate}
	\item Ejecutar el pipeline completo sobre la imagen asignada.
	\item Realizar la calibración de escala interactiva (notebook).
	\item Medir al menos 15 características dimensionales.
	\item Diagnosticar el modo de fractura predominante.
	\item Completar la evaluación de riesgo
	      (Tabla~\ref{tab:risk_template}).
	\item Comparar resultados con al menos otro equipo.
\end{enumerate}

\begin{table}[H]
	\centering
	\caption{Plantilla de evaluación de riesgo (a completar por cada equipo).}
	\label{tab:risk_template}
	\small
	\begin{tabularx}{\textwidth}{l c X}
		\toprule
		\textbf{Factor}                   & \textbf{Valor (1--5)} & \textbf{Justificación}                                           \\
		\midrule
		Probabilidad ($P$)                &                       & (MuyBajo=1, Bajo=2, Medio=3, Alto=4, MuyAlto=5)                  \\
		Severidad ($S$)                   &                       & (Insignificante=1, Menor=2, Moderado=3, Mayor=4, Catastrófico=5) \\
		\midrule
		Riesgo $R = P \times S$           &                       &                                                                  \\
		Nivel de riesgo                   &                       & Bajo ($<5$) / Medio (5--12) / Alto ($>12$)                       \\
		\midrule
		Controles requeridos              & {--}                  & (Listar: EPP, ingeniería, administrativos)                       \\
		Posición en la matriz (fila, col) & {--}                  & (Marcar en Figura~\ref{fig:risk_matrix})                         \\
		\bottomrule
	\end{tabularx}
\end{table}

% ---- EQUIPO 1 ----
\subsection{Equipo~1: Análisis dimensional de estrías de fatiga}
\label{sec:eq1}

\noindent\fbox{\begin{minipage}{0.96\textwidth}
	\textbf{Imagen asignada:} \texttt{fatigue3.png}

	\textbf{Contexto:} Se sospecha que el componente falló por fatiga de alto
	ciclo.  Las estrías de fatiga son evidencia directa del avance de la
	grieta.  El espaciado entre estrías permite estimar $da/dN$.

	\textbf{Tareas:}
	\begin{enumerate}
		\item Ejecutar el pipeline completo y documentar los 6 resultados.
		\item Calibrar la escala con la barra del SEM.
		\item Medir el espaciado entre 15--20 estrías consecutivas.
		\item Calcular el espaciado promedio $\bar{s}$ y la desviación
		      estándar $\sigma_s$.
		\item Estimar el rango del factor de intensidad de tensiones $\Delta K$
		      usando la relación empírica $da/dN \approx \bar{s}$ y la ley
		      de Paris (ec.~\ref{eq:paris}) con $C = 6{,}9 \times 10^{-12}$
		      y $m = 3$ (unidades SI).
		\item Comparar el valor de $\Delta K$ con el umbral de propagación
		      $\Delta K_\mathrm{th} \approx \SI{5}{\MPa\sqrt{\m}}$ para
		      aleaciones de aluminio aeronáuticas.
	\end{enumerate}
\end{minipage}}

\textbf{Evaluación de riesgo:} Complete la Tabla~\ref{tab:risk_template}
considerando los riesgos de operación del SEM y preparación metalográfica.

\subsubsection*{Problemas Adicionales -- Equipo~1}

\begin{enumerate}
	\item[\textbf{P1.1}] Si el error en la calibración de la barra de escala
	      es de $\pm$5~píxeles y la barra mide 200~píxeles, ¿cuál es el error
	      porcentual en las mediciones de espaciado? ¿Cómo se propaga a
	      $\Delta K$?

	\item[\textbf{P1.2}] Compare el número de regiones detectadas por Otsu
	      usando \texttt{clip\_limit}~$= 0{,}01$ vs.\ $0{,}05$.  ¿El CLAHE
	      mejora la segmentación?

	\item[\textbf{P1.3}] Proponga un criterio cuantitativo (basado en los
	      descriptores de forma) para distinguir automáticamente estrías de
	      fatiga de hoyuelos.  Sugerencia: considere la excentricidad y la
	      relación área/perímetro.
\end{enumerate}

% ---- EQUIPO 2 ----
\subsection{Equipo~2: Comparación de modos de fractura}
\label{sec:eq2}

\noindent\fbox{\begin{minipage}{0.96\textwidth}
	\textbf{Imágenes asignadas:} \texttt{fatigue4.png} y \texttt{fatigue7.png}

	\textbf{Contexto:} Dos zonas del mismo componente fracturado muestran
	morfologías diferentes.  Se requiere determinar si representan fractura
	dúctil, frágil o mixta.

	\textbf{Tareas:}
	\begin{enumerate}
		\item Procesar ambas imágenes con el pipeline.
		\item Para cada imagen: identificar si predominan hoyuelos (dúctil)
		      o facetas planas (frágil).
		\item Calibrar y medir el diámetro de al menos 15 hoyuelos (si los hay)
		      o el tamaño de 15 facetas.
		\item Construir un histograma con las mediciones y calcular media y
		      desviación estándar.
		\item Comparar los descriptores de forma (excentricidad media, área
		      media) entre ambas imágenes.
		\item Emitir un diagnóstico del modo de fractura para cada zona.
	\end{enumerate}
\end{minipage}}

\textbf{Evaluación de riesgo:} Complete la Tabla~\ref{tab:risk_template}.

\subsubsection*{Problemas Adicionales -- Equipo~2}

\begin{enumerate}
	\item[\textbf{P2.1}] Si una zona muestra hoyuelos equiaxiales y otra
	      hoyuelos alargados, ¿qué puede inferir sobre el estado de tensiones
	      en cada zona?

	\item[\textbf{P2.2}] Genere las imágenes Canny con $\sigma = 1{,}0$ y
	      $\sigma = 3{,}0$.  ¿Cómo afecta $\sigma$ a la detección de bordes
	      finos (estrías) vs.\ bordes gruesos (facetas)?

	\item[\textbf{P2.3}] Proponga una modificación al pipeline para detectar
	      fractura intergranular (límites de grano).  ¿Qué operaciones
	      adicionales serían necesarias?
\end{enumerate}

% ---- EQUIPO 3 ----
\subsection{Equipo~3: Efecto de la resolución y magnificación}
\label{sec:eq3}

\noindent\fbox{\begin{minipage}{0.96\textwidth}
	\textbf{Imágenes asignadas:} \texttt{fatigue5.png} (alta resolución) y
	\texttt{fatigue4.png} (baja resolución)

	\textbf{Contexto:} La magnificación afecta la cantidad de detalle
	resuelto y, por tanto, las métricas cuantitativas.  Se debe cuantificar
	este efecto.

	\textbf{Tareas:}
	\begin{enumerate}
		\item Procesar ambas imágenes con parámetros idénticos.
		\item Comparar el número de regiones detectadas ($N$).
		\item Comparar la distribución de áreas de las regiones (histograma).
		\item Calcular la relación $\mathrm{px}/\si{\micro\m}$ para cada
		      imagen y discutir cómo la resolución afecta la precisión de las
		      mediciones.
		\item Discutir si es válido comparar mediciones absolutas (en
		      \si{\micro\m}) entre imágenes de distinta magnificación.
	\end{enumerate}
\end{minipage}}

\textbf{Evaluación de riesgo:} Complete la Tabla~\ref{tab:risk_template}.

\subsubsection*{Problemas Adicionales -- Equipo~3}

\begin{enumerate}
	\item[\textbf{P3.1}] Si \texttt{fatigue5.png} tiene una resolución de
	      $0{,}1$~\si{\micro\m}/px y \texttt{fatigue4.png} de
	      $0{,}5$~\si{\micro\m}/px, ¿cuál es la menor característica medible
	      en cada caso?  (Criterio: 3 píxeles mínimo.)

	\item[\textbf{P3.2}] Redimensione \texttt{fatigue5.png} al 25\% de su
	      tamaño original con \texttt{cv2.resize} y reprocese.  ¿Cuántas
	      regiones se detectan?  ¿Cuánto cambia respecto al original?

	\item[\textbf{P3.3}] Proponga un protocolo experimental (magnificaciones,
	      número de imágenes, zonas) para obtener estadísticas representativas
	      de una superficie de fractura completa.
\end{enumerate}

% ---- EQUIPO 4 ----
\subsection{Equipo~4: Optimización de parámetros del pipeline}
\label{sec:eq4}

\noindent\fbox{\begin{minipage}{0.96\textwidth}
	\textbf{Imagen asignada:} \texttt{fatigue6.png} (la más compleja)

	\textbf{Contexto:} Los parámetros del pipeline (clip\_limit, $\sigma$
	de Canny, tamaño de disco morfológico) afectan los resultados.  Se
	busca optimizar estos parámetros para maximizar la calidad de la
	segmentación.

	\textbf{Tareas:}
	\begin{enumerate}
		\item Procesar con parámetros por defecto (clip\_limit\,=\,0,03;
		      $\sigma = 2{,}0$; disk\,=\,2).
		\item Variar \texttt{clip\_limit} $\in \{0{,}01;\; 0{,}03;\; 0{,}05;\;
		      0{,}10\}$ y registrar $N_\mathrm{regiones}$.
		\item Variar $\sigma_\mathrm{Canny} \in \{1{,}0;\; 2{,}0;\; 3{,}0;\;
		      5{,}0\}$ y registrar $N_\mathrm{regiones}$.
		\item Variar disk $\in \{1, 2, 3, 5\}$ y registrar
		      $N_\mathrm{regiones}$.
		\item Generar una tabla resumen y gráficas de sensibilidad.
		\item Proponer la combinación óptima justificada.
	\end{enumerate}
\end{minipage}}

\textbf{Evaluación de riesgo:} Complete la Tabla~\ref{tab:risk_template}.

\subsubsection*{Problemas Adicionales -- Equipo~4}

\begin{enumerate}
	\item[\textbf{P4.1}] ¿Existe un valor de disk suficientemente grande que
	      elimine \emph{todas} las regiones?  Calcúlelo teóricamente usando la
	      distribución de áreas.

	\item[\textbf{P4.2}] Implemente la operación de \emph{closing} (cierre)
	      después de opening.  ¿Cómo cambia el número de regiones y la forma
	      de los objetos?

	\item[\textbf{P4.3}] Compare los resultados de Otsu con un umbral manual
	      fijo (e.g., $t = 128$).  ¿En qué condiciones Otsu falla?
\end{enumerate}

% =============================================================================
\section{Evaluación de Riesgos}
\label{sec:riesgos}

\subsection{Matriz de Riesgos}

\begin{figure}[H]
	\centering
	\begin{tikzpicture}[scale=0.85]
		\foreach \x in {0,...,4} {
			\foreach \y in {0,...,4} {
				\pgfmathsetmacro{\risk}{(\x+1)*(\y+1)}
				\pgfmathparse{\risk < 5 ? "aerogreen!40" :
					(\risk < 13 ? "aeroyellow!60" : "aerored!50")}
				\edef\fillcol{\pgfmathresult}
				\fill[\fillcol] (\x,\y) rectangle (\x+1,\y+1);
				\draw[gray] (\x,\y) rectangle (\x+1,\y+1);
				\node at (\x+0.5,\y+0.5)
					{\pgfmathprintnumber[fixed,precision=0]{\risk}};
			}
		}
		\foreach \x/\l in {0/MuyBajo,1/Bajo,2/Medio,3/Alto,4/MuyAlto} {
			\node[font=\tiny,rotate=45,anchor=west] at (\x+0.5,-0.1) {\l};
		}
		\foreach \y/\l in {0/Insignificante,1/Menor,2/Moderado,3/Mayor,
			4/Catastrófico} {
			\node[font=\tiny,anchor=east] at (-0.1,\y+0.5) {\l};
		}
		\node[font=\small\bfseries,rotate=90] at (-1.5,2.5) {Severidad};
		\node[font=\small\bfseries] at (2.5,-1.2) {Probabilidad};
		% Scenario markers
		\fill[aerored,opacity=0.8] (1.5,3.5) circle (0.25);
		\node[white,font=\tiny\bfseries] at (1.5,3.5) {SEM};
		\fill[blue!70,opacity=0.8] (1.5,2.5) circle (0.25);
		\node[white,font=\tiny\bfseries] at (1.5,2.5) {Prep.};
		\fill[orange,opacity=0.8] (1.5,2.5) circle (0.25);
		\node[white,font=\tiny\bfseries] at (1.5,2.5) {Reac.};
		\fill[aerogreen!70,opacity=0.8] (0.5,0.5) circle (0.25);
		\node[white,font=\tiny\bfseries] at (0.5,0.5) {PC};
		% Legend
		\fill[aerogreen!40] (5.5,3.5) rectangle (6.0,4.0);
		\node[right,font=\tiny] at (6.0,3.75) {Bajo ($<5$)};
		\fill[aeroyellow!60] (5.5,2.8) rectangle (6.0,3.3);
		\node[right,font=\tiny] at (6.0,3.05) {Medio (5--12)};
		\fill[aerored!50] (5.5,2.1) rectangle (6.0,2.6);
		\node[right,font=\tiny] at (6.0,2.35) {Alto ($>12$)};
	\end{tikzpicture}
	\caption{Matriz de riesgos $5\times5$ para operaciones de análisis de
		imagen con SEM\@.  SEM: zona amarilla; preparación metalográfica y
		reactivos: zona amarilla; procesamiento digital: zona verde.}
	\label{fig:risk_matrix}
\end{figure}

\subsection{Riesgos identificados}

\begin{table}[H]
	\centering
	\caption{Identificación y control de riesgos en la práctica de análisis
		de imagen.}
	\label{tab:riesgos}
	\small
	\begin{tabularx}{\textwidth}{l c c c X}
		\toprule
		\textbf{Riesgo}       & \textbf{P} & \textbf{S} & \textbf{R}
		                      & \textbf{Medidas de control}                       \\
		\midrule
		Alta tensión del SEM
		                      & 2          & 4          & 8
		                      & Interlock de cámara, formación obligatoria,
		puesta a tierra, nunca abrir durante operación.                           \\
		Preparación muestras
		                      & 2          & 3          & 6
		                      & Gafas de protección, guantes, fijación correcta
		de la muestra en la pulidora.                                              \\
		Reactivos metalográficos
		                      & 2          & 3          & 6
		                      & Campana de extracción, guantes de nitrilo,
		hojas MSDS disponibles, lavaojos.                                         \\
		Ergonomía digital
		                      & 1          & 1          & 1
		                      & Pausas cada 45~min, iluminación adecuada,
		posición ergonómica del monitor.                                           \\
		\bottomrule
	\end{tabularx}
\end{table}

\subsection{Jerarquía de controles}

\begin{enumerate}
	\item \textbf{Eliminación/Sustitución}: cuando sea posible, usar réplicas
	      poliméricas de la superficie de fractura para evitar manipulación
	      directa de piezas metálicas afiladas.
	\item \textbf{Controles de ingeniería}: interlock de puerta del SEM,
	      campanas de extracción, protectores en pulidoras.
	\item \textbf{Controles administrativos}: SOPs para operación del SEM,
	      formación obligatoria, supervisión del responsable de laboratorio,
	      registro de uso del equipo.
	\item \textbf{EPP}: gafas de seguridad para preparación de muestras,
	      guantes de nitrilo para reactivos, bata de laboratorio.
\end{enumerate}

% =============================================================================
\section{Herramientas de Software}
\label{sec:software}

Carpeta \texttt{Practicas/P05\_Analisis\_Imagen/}:

\subsection{Script principal \texttt{P05\_Analisis\_Imagen.py}}

El script \texttt{src/P05\_Analisis\_Imagen.py} (\textasciitilde258 líneas)
ejecuta el pipeline completo de análisis en modo \emph{batch} sobre todas
las imágenes encontradas en \texttt{files/p5/}:

\begin{table}[H]
	\centering
	\caption{Funciones principales del script \texttt{P05\_Analisis\_Imagen.py}.}
	\label{tab:funciones}
	\small
	\begin{tabularx}{\textwidth}{l l X}
		\toprule
		\textbf{Función}        & \textbf{Firma}          & \textbf{Comportamiento}           \\
		\midrule
		\texttt{cargar\_imagen\_gris}
		                        & \texttt{(path)}
		& Carga PNG/JPG en gris, normaliza a float $[0,1]$.                                     \\
		\texttt{procesar\_clahe}
		                        & \texttt{(img, clip)}
		& Aplica CLAHE con \texttt{clip\_limit} dado.                                           \\
		\texttt{procesar\_sobel}
		                        & \texttt{(img)}
		& Filtro Sobel $\to$ mapa de gradientes.                                                 \\
		\texttt{procesar\_canny}
		                        & \texttt{(img, $\sigma$)}
		& Detector Canny $\to$ bordes binarios.                                                  \\
		\texttt{procesar\_otsu}
		                        & \texttt{(img)}
		& Umbral de Otsu $\to$ máscara binaria.                                                  \\
		\texttt{procesar\_morfologia}
		                        & \texttt{(img, op, disk)}
		& Opening/closing/erosion/dilation con disco indicado.                                   \\
		\texttt{procesar\_etiquetado}
		                        & \texttt{(img\_bin)}
		& Etiquetado de regiones $\to$ propiedades (área, perímetro,
		excentricidad, eje mayor).                                                                \\
		\texttt{\_save\_individual}
		                        & \texttt{(img, cmap, title, file)}
		& Crea, guarda y cierra una figura individual para cada paso.                            \\
		\texttt{run\_full\_analysis}
		                        & \texttt{(img, base)}
		& Ejecuta los 6 pasos + figura compuesta + etiquetado.                                  \\
		\bottomrule
	\end{tabularx}
\end{table}

\begin{lstlisting}[language=Python,caption={Uso del script desde la línea de
  comandos.}]
# Activar entorno virtual
.venv\Scripts\Activate.ps1

# Ejecutar desde la raiz del repositorio
python Practicas\P05_Analisis_Imagen\src\P05_Analisis_Imagen.py

# Resultados en: Practicas/P05_Analisis_Imagen/data/images/
# 30 archivos PNG (6 por imagen x 5 imagenes)
\end{lstlisting}

\subsection{Notebook Jupyter}

\texttt{notebooks/P05\_Analisis\_Imagen.ipynb} reproduce el análisis de
forma interactiva e incluye funcionalidad adicional:

\begin{itemize}
	\item \textbf{Widget de selección}: menú desplegable para elegir la
	      imagen a analizar.
	\item \textbf{Calibración interactiva}: mediante la clase
	      \texttt{CalculadoraDistancia} (clics en la barra de escala del SEM).
	\item \textbf{Medición interactiva}: clics sobre la micrografía para
	      medir distancias en \si{\micro\m} utilizando la escala calibrada.
	\item \textbf{Análisis automático}: al seleccionar una imagen, se ejecuta
	      el pipeline completo y se muestran los resultados \emph{inline}.
	\item \textbf{Guardado idéntico}: genera los mismos 6 archivos por imagen
	      que el script \texttt{.py}.
\end{itemize}

\subsection{Datos de salida}

\texttt{data/images/} almacena los resultados procesados:
\begin{itemize}
	\item 5 $\times$ \texttt{\_00\_resumen.png}: figura compuesta
	      $3\times2$.
	\item 5 $\times$ \texttt{\_01\_clahe.png} a
	      \texttt{\_05\_opening.png}: etapas individuales.
	\item Total: 30 archivos PNG.
\end{itemize}

% =============================================================================
\section{Procedimiento Experimental}
\label{sec:procedimiento}

\begin{enumerate}
	\item \textbf{Preparación de la muestra} (si se realiza en laboratorio):
	\begin{enumerate}
		\item Seleccionar el espécimen de fractura.
		\item Montar en porta-muestras del SEM con cinta conductora de
		      carbono.
		\item Recubrir con oro-paladio si el material no es conductor
		      (sputtering, $\sim\!\SI{10}{\nm}$).
	\end{enumerate}

	\item \textbf{Adquisición SEM}:
	\begin{enumerate}
		\item Insertar la muestra en la cámara y evacuar.
		\item Seleccionar EHT (típicamente 5--20~kV), WD y detector (SE
		      para topografía).
		\item Adquirir micrografías a distintas magnificaciones.
		\item Verificar que la barra de escala esté visible y sea legible.
	\end{enumerate}

	\item \textbf{Análisis de imagen digital}:
	\begin{enumerate}
		\item Ejecutar el script \texttt{P05\_Analisis\_Imagen.py} o abrir
		      el notebook.
		\item Verificar que las 5 imágenes se procesan correctamente
		      (30 archivos de salida).
		\item En el notebook: calibrar la escala con la barra del SEM.
		\item Realizar las mediciones dimensionales requeridas por el caso
		      del equipo.
		\item Documentar capturas de pantalla de cada etapa del pipeline.
	\end{enumerate}

	\item \textbf{Diagnóstico fractográfico}:
	\begin{enumerate}
		\item Identificar el modo de fractura predominante.
		\item Correlacionar las observaciones cualitativas con los datos
		      cuantitativos (número de regiones, dimensiones de
		      características).
		\item Emitir un diagnóstico preliminar.
	\end{enumerate}
\end{enumerate}

% =============================================================================
\section{Cuestionario}
\label{sec:cuestionario}

\begin{enumerate}
	\item Explique por qué la calibración precisa de la escala es fundamental
	      para cualquier medición cuantitativa.  ¿Qué error porcentual se
	      introduce si la selección de la barra tiene un error de
	      $\pm$5~píxeles sobre una barra de 200~píxeles?

	\item Describa la diferencia entre el filtro Sobel y el detector de
	      Canny.  ¿En qué situaciones cada uno es más apropiado para el
	      análisis fractográfico?

	\item El método de Otsu minimiza la varianza intra-clase.  Escriba la
	      expresión de la varianza inter-clase $\sigma_B^2(t)$ y explique por
	      qué maximizarla es equivalente a minimizar la varianza intra-clase.

	\item ¿Por qué el SEM requiere que la muestra sea conductora o esté
	      recubierta?  ¿Qué artefactos aparecen en la imagen si la muestra
	      no conduce?

	\item La operación de apertura (opening) elimina objetos pequeños.
	      Demuestre matemáticamente que $I \circ SE \subseteq I$ (la apertura
	      nunca agrega píxeles blancos).

	\item Si se observan hoyuelos muy pequeños y equiaxiales, ¿qué infiere
	      sobre la ductilidad del material?  Compare con hoyuelos grandes
	      y alargados.

	\item ¿Por qué es importante analizar la fractura a diferentes
	      magnificaciones?  ¿Qué información proporciona cada escala?

	\item El voltaje de aceleración (EHT) del SEM afecta la penetración
	      del haz.  ¿Cómo influye en el contraste de la imagen resultante?
	      ¿Es preferible alto o bajo EHT para fractografía superficial?

	\item Si sospecha fatiga pero no identifica estrías claras,
	      ¿qué otras características buscaría?  ¿Qué imágenes adicionales
	      solicitaría al operador del SEM?

	\item Proponga un método cuantitativo, basado en los descriptores de
	      forma de las regiones etiquetadas, para clasificar automáticamente
	      una superficie como ``predominantemente dúctil'' o
	      ``predominantemente frágil''.
\end{enumerate}

% =============================================================================
\section{Formato del Informe}
\label{sec:informe}

Formato AIAA con las siguientes secciones obligatorias:

\begin{enumerate}
	\item \textbf{Resumen}: objetivo, metodología (pipeline de 6 etapas,
	      calibración, mediciones), resultados principales (modo de fractura
	      diagnosticado, dimensiones medidas).
	\item \textbf{Introducción}: importancia de la fractografía SEM,
	      objetivos de la práctica.
	\item \textbf{Fundamento Teórico}: principios SEM, mecanismos de
	      fractura, algoritmos de procesamiento.
	\item \textbf{Metodología}: descripción del pipeline, parámetros
	      utilizados, procedimiento de calibración y medición.
	\item \textbf{Resultados}: capturas de las 6 etapas del pipeline,
	      tabla de regiones, mediciones dimensionales, histogramas.
	\item \textbf{Discusión}: diagnóstico fractográfico, respuestas al
	      cuestionario, efecto de los parámetros del pipeline.
	\item \textbf{Conclusiones}: resumen de hallazgos, validez del
	      diagnóstico, recomendaciones.
	\item \textbf{Referencias}: formato IEEE.
	\item \textbf{Apéndice~A}: imágenes generadas por el script/notebook.
\end{enumerate}

% =============================================================================
\section{Conclusiones}
\label{sec:conclusiones}

\begin{enumerate}
	\item El pipeline de 6 etapas (carga, CLAHE, Sobel/Canny, Otsu,
	      opening, etiquetado) permite procesar sistemáticamente micrografías
	      SEM de superficies de fractura, extrayendo información tanto
	      cualitativa como cuantitativa.

	\item La ecualización adaptativa CLAHE es esencial para micrografías
	      con iluminación desigual, mejorando el contraste local sin
	      amplificar el ruido globalmente (controlado por
	      \texttt{clip\_limit}).

	\item El método de Otsu proporciona un umbral de segmentación óptimo
	      y automático, evitando la selección manual que introduce
	      subjetividad.  Su validez depende de que el histograma sea
	      bimodal.

	\item Las operaciones morfológicas (apertura con disco) son
	      indispensables para eliminar ruido residual y obtener regiones
	      que representen características reales de la superficie.

	\item El número de regiones detectadas varía significativamente con
	      la resolución de la imagen (153 a 2897 regiones), lo que
	      demuestra la importancia de reportar siempre las condiciones de
	      magnificación y los parámetros del análisis.

	\item La calibración mediante la barra de escala del SEM es el paso
	      más crítico para la fiabilidad de las mediciones dimensionales;
	      un error de 5~píxeles puede introducir un error del 5\% en las
	      mediciones.

	\item Las herramientas de software (\texttt{P05\_Analisis\_Imagen.py}
	      y notebook) automatizan el flujo completo --- desde la detección
	      de imágenes hasta la generación de 30 archivos de resultados ---
	      garantizando reproducibilidad y trazabilidad del análisis.
\end{enumerate}

% =============================================================================
\printbibliography[title={Referencias}]

% =============================================================================
\appendix

\section{Referencia del Pipeline de Procesamiento}
\label{sec:appendix_pipeline}

\begin{table}[H]
	\centering
	\caption{Resumen de parámetros del pipeline y sus efectos.}
	\label{tab:parametros_pipeline}
	\small
	\begin{tabularx}{\textwidth}{l l l X}
		\toprule
		\textbf{Etapa} & \textbf{Parámetro} & \textbf{Valor}
		& \textbf{Efecto al aumentar}                                           \\
		\midrule
		CLAHE          & \texttt{clip\_limit}     & 0,03
		& Mayor amplificación del contraste local; riesgo de ruido.            \\
		Canny          & $\sigma$                  & 2,0
		& Mayor suavizado previo $\to$ menos bordes detectados, más gruesos.   \\
		Morfología     & disk size                  & 2
		& Elimina objetos más grandes; reduce $N_\mathrm{regiones}$.           \\
		Otsu           & ---                        & Automático
		& Umbral óptimo $t^*$; no tiene parámetro de usuario.                  \\
		\bottomrule
	\end{tabularx}
\end{table}

\section{Referencia Rápida de Funciones Python}
\label{sec:appendix_python}

\begin{table}[H]
	\centering
	\caption{Referencia rápida: funciones de \texttt{P05\_Analisis\_Imagen.py}.}
	\label{tab:ref_python}
	\small
	\begin{tabularx}{\textwidth}{l X}
		\toprule
		\textbf{Función} & \textbf{Uso rápido}                              \\
		\midrule
		\texttt{cargar\_imagen\_gris(path)}
		& $\rightarrow$ \texttt{ndarray} float $[0,1]$                       \\
		\texttt{procesar\_clahe(img, 0.03)}
		& $\rightarrow$ Imagen con contraste mejorado                        \\
		\texttt{procesar\_sobel(img)}
		& $\rightarrow$ Mapa de gradientes float                             \\
		\texttt{procesar\_canny(img, 2.0)}
		& $\rightarrow$ Imagen binaria de bordes                             \\
		\texttt{procesar\_otsu(img)}
		& $\rightarrow$ Máscara binaria (bool)                               \\
		\texttt{procesar\_morfologia(img, 'opening', 2)}
		& $\rightarrow$ Máscara refinada                                     \\
		\texttt{procesar\_etiquetado(img\_bin)}
		& $\rightarrow$ \texttt{(label\_image, props\_table)}                \\
		\texttt{run\_full\_analysis(img, 'fatigue3')}
		& Ejecuta pipeline + guarda 6 PNG + imprime $N$ regiones            \\
		\bottomrule
	\end{tabularx}
\end{table}

\section{Formulario Matemático}
\label{sec:appendix_formulas}

\begin{table}[H]
	\centering
	\caption{Resumen de ecuaciones principales.}
	\small
	\begin{tabularx}{\textwidth}{l X l}
		\toprule
		\textbf{Ecuación}       & \textbf{Expresión}
		& \textbf{Ref.}                                                   \\
		\midrule
		Sobel                   & $|\nabla I| = \sqrt{G_x^2 + G_y^2}$
		& \eqref{eq:sobel_mag}                                            \\
		Gaussiano (Canny)       & $G_\sigma = \frac{1}{2\pi\sigma^2}\exp(-\tfrac{r^2}{2\sigma^2})$
		& \eqref{eq:gauss_kernel}                                         \\
		Otsu                    & $t^* = \arg\max\;\omega_0\omega_1(\mu_0-\mu_1)^2$
		& \eqref{eq:otsu_opt}                                             \\
		Opening                 & $I \circ SE = (I\ominus SE)\oplus SE$
		& \eqref{eq:opening}                                              \\
		Calibración             & $s = d_\mathrm{px}/d_\mathrm{real}$
		& \eqref{eq:escala}                                               \\
		Error de escala         & $\varepsilon_s = 2n/d_\mathrm{px}\times100\%$
		& \eqref{eq:error_escala}                                         \\
		Ley de Paris            & $da/dN = C(\Delta K)^m$
		& \eqref{eq:paris}                                                \\
		\bottomrule
	\end{tabularx}
\end{table}

\end{document}
